\section{Quadratic charts and local detection}\label{sec:charts}

This section establishes the scalar calculus used in the subsequent
projection constructions.  The distinction between a chart's complexity
and its real oscillation parameters is essential: the latter will never
be bounded.  All constants in this section are independent of those
parameters and of the interval on which an atom is tested.

We first combine a scale-wise Weyl alternative and Fourier expansion
with an algebraic separation identity along the progression, obtaining
separated moments at one scale.  We then transfer the discrete
degree-two inverse theorem to bounded continuous witnesses and use
Hahn--Banach to decompose bounded inputs.  Applying that decomposition
to two opposing inputs reduces local detection to the separated moments
already established.

\subsection{Atoms and quantitative conventions}

Fix once and for all the Gevrey order used for the smooth cutoffs.
Enlarging its fixed constants does not change any of the estimates below.
For a rational number, its height is the maximum of the absolute
numerator and the positive denominator in reduced form.  The height of
a rational vector or matrix is the maximum of the heights of its entries.
We use Euclidean norms for real vectors and the associated matrix norms;
changing between these norms and coordinate maximum norms only changes
the quantitative constants.

\begin{definition}[Quadratic chart atom]\label{def:quadratic-atom}
A quadratic chart atom is a function
\begin{equation*}
 \varphi(y)=e(\lambda y^2+\beta y)
 \sum_{n\in\Z^d}\chi(\theta y-n)
                 e\bigl((\theta y)^T C(\theta y-n)\bigr).
 \tag{1}\label{eq:src-1}
\end{equation*}
Here \(\theta\in\R^d\), \(\lambda,\beta\in\R\), \(C\) is a rational
\(d\)-by-\(d\) matrix, and \(\chi\) is smooth and compactly supported.
The case \(d=0\) includes ordinary quadratic phases.  An overall
unimodular constant factor is allowed.  Every non-unimodular scalar
multiplying an atom is either absorbed into \(\chi\), and is then
subject to its stated complexity bounds, or is written as an external
coefficient.

An atom has complexity at most \(H\), where \(H\ge e^2\), if
\(d\le\log H\), the height of \(C\) is at most \(H\),
\(\supp\chi\subset[-H,H]^d\), and, for every multi-index \(\alpha\),
\[
 \|\partial^\alpha\chi\|_\infty
 \le H^{c_0(|\alpha|+1)}
          (|\alpha|+1)^{c_0(|\alpha|+1)} .
\]
The constant \(c_0\) is fixed throughout.  The parameters
\(\lambda,\beta,\theta\) are unrestricted.
A bounded atom means an atom normalized to have supremum norm at most
one.
\end{definition}

A quantity described below as quasipolynomial in \(H_1,\ldots,H_m\) is
bounded by
\[
 \exp\!\left(C\left(1+\sum_{\nu=1}^m\log H_\nu\right)^C\right),
 \qquad H_\nu\ge2,
\]
with an absolute \(C\); a fixed number of compositions of such bounds
is again of this form.  When \(q>2\) occurs, the constant multiplying
the logarithms may depend on \(q\), but the exponent \(C\) can be chosen
independently of \(q\).  Constants attached to a fixed requested
Fourier-tail power may also depend on that power.

There are at most \((2H+3)^d\) nonzero summands at a point in
\eqref{eq:src-1}.  Consequently its supremum norm is quasipolynomial in
\(H\).  Dividing an atom by such a bound normalizes it; the factor is
then included in the coefficient of the expansion in which it occurs.

\begin{lemma}[Elementary chart operations]\label{lem:chart-operations}
Complex conjugation, a product of a bounded number of atoms, and
translation or affine rescaling of the real argument preserve the class
\eqref{eq:src-1}, with a quasipolynomial change of complexity.  The same
statement holds for the multivariable charts introduced below.
Only the unrestricted real oscillation parameters can depend on the
size of the translation or rescaling.
\end{lemma}

\begin{proof}
For conjugation replace \(C,\lambda,\beta\) by their negatives and
conjugate \(\chi\).  For products use the direct sum of the horizontal
coordinates and the block diagonal sum of the matrices.  The product
rule for derivatives gives the required Gevrey bounds.

It remains to check a translated horizontal argument.  Let that argument
be \(\theta y+a\).  Choose a positive integer \(P\) such that every
entry of \(PC\) is integral.  The product of the denominators of the
\(d^2\) entries is an admissible choice, so
\(P\le H^{d^2}\).  Write \(a=m_0+a_0\), with
\(m_0\in P\Z^d\) and \(\|a_0\|_\infty\le P\), and replace the summation
index by \(n+m_0\).  With \(w=\theta y-n\), the new phase is
\[
 (\theta y+m_0+a_0)^TC(w+a_0)
 =(\theta y)^TCw
    +(\theta y)^TCa_0+m_0^TC\theta y
    +a_0^TC(w+a_0)+m_0^TCa_0
       \pmod{\Z}.
\]
The two middle terms linear in \(\theta y\) change the unrestricted
linear phase.  The last term is constant.  The term
\(a_0^TC(w+a_0)\) can be included in the smooth amplitude
\(\chi(w+a_0)\).  The support and derivatives of that amplitude have
quasipolynomial bounds because \(a_0\), unlike \(m_0\), is bounded by
\(P\).  Translating an ordinary quadratic phase produces only a
quadratic phase, a linear phase, and a constant.  Rescaling simply
rescales \(\theta,\lambda,\beta\).
All these calculations hold with \(\theta y\) replaced by \(Vz\).
\end{proof}

\subsection{A quantitative Weyl alternative}

For \(z\in\R^k\), where \(1\le k\le3\), consider
\[
 Q(z)=e(z^T\Lambda z+b^Tz)
       \sum_{n\in\Z^d}\chi(Vz-n)
                   e\bigl((Vz)^TC(Vz-n)\bigr),
 \qquad A=C-C^T.
\]
The matrix \(\Lambda\) is real and symmetric and \(V\) is real.
The complexity conditions apply to \(d,C,\chi\) as in
Definition~\ref{def:quadratic-atom}.  We write
\[
 C_{\mathrm{sy}}=\frac{C+C^T}{2},\qquad
 C_{\mathrm{sk}}=\frac{C-C^T}{2},\qquad
 \operatorname{sym}B=\frac{B+B^T}{2}.
\]

\begin{lemma}[Scale-wise Weyl alternative]\label{lem:weyl}
Let \(\vartheta\) be supported in a fixed box in \(\R^k\), with
supremum norm and Lipschitz norm at most \(H\).
Suppose \(0<\delta<1/2\) and
\[
 \sup_{\xi\in\R^k}
 \left|\int_{\R^k}\vartheta(z)Q(z)e(-\xi^Tz)\,dz\right|
 \ge\delta .
\]
There is a rational subspace \(U\subset\R^d\), with orthogonal
projections \(P_f\) onto \(U\) and \(P_s=1-P_f\), such that
\begin{equation*}
 \begin{gathered}
 A|_{U\times U}=0,\qquad \|V_s\|\le H',\qquad
 \|\Lambda_*\|\le H',\\
 V_f=P_fV,\quad V_s=P_sV,\qquad
 \Lambda_*=\Lambda+\tfrac12V^TC_{\mathrm{sy}}V+
                     \operatorname{sym}(V_f^TC_{\mathrm{sk}}V_s).
 \end{gathered}
 \tag{2}\label{eq:src-2}
\end{equation*}
The heights of \(P_f,P_s\) are at most \(H'\), where
\[
 H'\le
 \exp\!\left(C(1+\log H+\log(1/\delta))^C\right).
\]
\end{lemma}

\begin{proof}
We first show that every sufficiently small return has a bounded
frequency increment.  The lattice directions generated by these returns
form a rational isotropic subspace with a bounded-height projection and
bounded slow horizontal matrix.  Returns in each coordinate direction
then bound the remaining effective quadratic matrix.
Let \(P\) clear the denominators of \(C\), as in
Lemma~\ref{lem:chart-operations}.  Put \(F=\vartheta Q\).
The support of \(F\), its supremum norm, its \(L^2\) norm, and
\(\int(1+\|z\|)|F(z)|\,dz\) are bounded by a quantity \(D\)
quasipolynomial in \(H\).
The value of \(D\) can be enlarged finitely many times below.

Suppose
\[
 Vh=m+d_h,\qquad m\in P\Z^d.
\]
Translate \(z\) by \(h\) and replace the lattice index by \(n+m\).
On the support of the translated amplitude, \(Vz-n\) has norm bounded
in terms of the chart complexity.  Expanding the phase therefore gives
\[
 F(z+h)=e(c(h,m))e(\Omega(h,m)^Tz)F(z)+E_{h,m}(z),
 \qquad
 \Omega(h,m)=(2\Lambda+V^TCV)h-V^TAm,
\]
where \(c(h,m)\) is constant in \(z\), and
\[
 \|E_{h,m}\|_{L^1}+\|E_{h,m}\|_\infty
       \le D(\|h\|+\|d_h\|)
\]
provided \(\|h\|+\|d_h\|\le1\).  To verify that the error contains no
unbounded coefficient, put \(u=Vz\), \(w=u-n\), and \(d=d_h\).
The nonconstant increment of the chart phase is
\[
 u^TCd+m^TCu+d^TCw.
\]
The first two terms, together with the ordinary quadratic increment,
are exactly \(\Omega(h,m)^Tz\); the term \(m^TCn\) discarded in this
calculation is integral.  The remaining \(d^TCw\) is bounded by
\(D\|d\|\).  Translating the amplitude and \(\vartheta\) costs
\(D(\|h\|+\|d\|)\).  This proves the claimed estimate.

The Fourier transform of \(F\) is \(D\)-Lipschitz.  Choose a point
\(\xi_0\) where its absolute value is at least \(3\delta/4\), and put
\(\rho=\delta/(100D)\).  Choose an integer
\[
 M\ge 100D^2\delta^{-2}\rho^{-k}.
\]
If necessary, increase \(D\) to cover the fixed volume constant in
Plancherel's theorem.  Choose \(a>0\) so small that
\[
 a\le\frac{\delta}{100DM},\qquad
 a\le\frac1{100(1+\|A\|)}.
\]
All of \(M,a^{-1},\rho^{-1}\) are quasipolynomial in \(H,\delta^{-1}\).
For a return with \(\|h\|+\|d_h\|\le a\), the same translation
calculation applies to \(\ell(h,m)\), \(1\le\ell\le M\).  It produces
Fourier values of absolute value at least \(\delta/2\) at
\(\xi_0-\ell\Omega(h,m)\).  If \(\|\Omega(h,m)\|\ge1\), the balls of
radius \(\rho\) about those points are disjoint, and on each such ball
the Fourier transform has magnitude at least \(\delta/3\).
Their contribution to its squared \(L^2\) norm exceeds \(D^2\), by the
choice of \(M\), contradicting Plancherel.  Hence
\[
 \|\Omega(h,m)\|<1
 \quad\hbox{whenever}\quad
 \|h\|+\|d_h\|\le a.
\]
For two returns satisfying this condition, direct antisymmetrization
gives
\[
 h'^T\Omega(h,m)-h^T\Omega(h',m')
       =d_{h'}^TAd_h-m'^TAm.
\]
The integer \(m'^TAm\) has absolute value less than one after decreasing
\(a\) by a fixed factor.  It is therefore zero.

We next find all directions generated by short returns.  For \(d\ge1\)
let
\[
 \mathcal K=P^{-1}V[-1,1]^k+[-1,1]^d.
\]
This is a symmetric convex body with nonempty interior.  Write
\(\lambda_1\le\cdots\le\lambda_d\) for its successive minima for
\(\Z^d\), and \(\mu_1\le\cdots\le\mu_d\) for the successive minima of
its polar body.  We need the elementary transference bounds
\[
 1\le\lambda_j\mu_{d+1-j}\le(d!)^2 .
\]
For the lower bound, independent lattice vectors in subspaces of
dimensions \(j\) and \(d+1-j\) cannot have all cross pairings zero.
One of their integer pairings has absolute value at least one.
For the upper bound, Minkowski's second theorem
\cite[Theorem~1.3]{Henk2002} gives
\[
 \prod_{j=1}^d\lambda_j\mu_j
 \le\frac{4^d}{\operatorname{vol}(\mathcal K)
                    \operatorname{vol}(\mathcal K^\circ)}
 \le(d!)^2.
\]
Here the last volume-product bound has a short direct proof.
Choose maximal-determinant columns \(v_1,\ldots,v_d\) in \(\mathcal K\).
Their cross-polytope is contained in \(\mathcal K\), and maximality
places \(\mathcal K\) in their coordinate parallelotope.
Taking polars gives an inscribed cross-polytope in
\(\mathcal K^\circ\) with reciprocal determinant.  The product of the
two cross-polytope volumes is \(4^d/(d!)^2\).
Since every opposite-minimum product is at least one, their product
upper bound proves each individual upper bound.

Set \(L=16(d!)^2\), and choose
\(b_{\max}\le a/(100P(d+k+1))\).
The \(d+1\) disjoint intervals
\[
 (b_{\max}L^{-j-1},\,b_{\max}L^{-j}],
 \qquad 0\le j\le d,
\]
cannot all contain a successive minimum.  Select one containing none,
write its left endpoint as \(b_0\), and put
\[
 U=\spanop(\Z^d\cap b_0\mathcal K).
\]
The reciprocal of \(b_0\) is quasipolynomial in the required parameters.
Every \(n\in\Z^d\cap b_0\mathcal K\) has a representation
\(n=P^{-1}Vh+w\), with \(\|h\|_\infty,\|w\|_\infty\le b_0\).
Thus \(m=Pn\), \(d_h=-Pw\) form a return satisfying the threshold \(a\).
The integrality argument above, applied to any two such \(n\), shows
that \(A\) vanishes on \(U\times U\).

Let \(r_0=\dim U\).  If \(r_0<d\), then
\(\lambda_{r_0+1}>Lb_0\); transference supplies \(d-r_0\) independent
integer vectors \(q_\nu\) in
\[
 \frac{(d!)^2}{Lb_0}\mathcal K^\circ .
\]
Each has zero pairing with \(\Z^d\cap b_0\mathcal K\), since that
pairing is integral and has absolute value at most \((d!)^2/L<1\).
They consequently span \(U^\perp\).  The inclusion
\([-1,1]^d\subset\mathcal K\) bounds their coordinates, while
\(P^{-1}V[-1,1]^k\subset\mathcal K\) bounds each \(q_\nu^TV\),
by quasipolynomial quantities.  If \(S\) has the \(q_\nu^T\) as rows,
\[
 P_s=S^T(SS^T)^{-1}S .
\]
Cramer's rule, and the positive integral determinant of \(SS^T\),
give quasipolynomial rational height and norm bounds for this matrix
and for \(P_sV\).  When \(r_0=d\), take \(P_s=0\); the same conclusions
hold.  Put \(P_f=1-P_s\).

Finally we bound the effective ordinary quadratic part.
For each coordinate vector \(e_j\in\R^k\), sample
\(P^{-1}V(te_j)\bmod\Z^d\) at equally spaced points
\(t\in[0,b_0]\).  Partition the horizontal torus into cubes of side
at most \(b_0/4\).  More than \((1+4/b_0)^d\) samples give two in the
same cube.  Their difference yields a return
\[
 h=t_je_j,\qquad b_0/M_0\le t_j\le b_0,\qquad
 \|d_h\|_\infty\le Pb_0/4,
\]
where \(M_0\le2(1+4/b_0)^d\).  Its lattice vector satisfies
\(P^{-1}m\in b_0\mathcal K\), so \(m\in U\).
In particular
\(m-V_fh=-P_fd_h\in U\).
Isotropy gives
\[
 V^TA(m-V_fh)=V_s^TA(m-V_fh),
\]
which is quasipolynomially bounded.  Direct matrix algebra gives
\[
 2\Lambda+V^TCV-V^TAV_f
       =2\Lambda_*+V_s^TC_{\mathrm{sk}}V_s .
\]
The return bound for \(\Omega\) therefore bounds
\(2\Lambda_*t_je_j\) quasipolynomially.  Since \(t_j^{-1}\) is also
quasipolynomial, this bounds every column of \(\Lambda_*\).
This proves \eqref{eq:src-2}.

For \(d=0\), the initial translation argument applies to every
sufficiently small \(h\), with \(\Omega=2\Lambda h\); take
\(h=ae_j/2\) and \(U=\{0\}\).  This also proves the assertion in that
case.
\end{proof}

\subsection{Fourier expansion after the Weyl splitting}

\begin{lemma}[Chart expansion and uniform tails]\label{lem:chart-expansion}
Suppose a rational orthogonal splitting satisfies
\eqref{eq:src-2}, with heights and the indicated sizes at most \(B\).
On a fixed box in \(\R^k\) there is an exact absolutely convergent
expansion
\begin{equation*}
 Q(z)=e(z^T\Lambda_*z+b^Tz)
                \sum_\ell a_\ell e(\sigma_\ell^TVz).
 \tag{3}\label{eq:src-3}
\end{equation*}
It consists of quasipolynomially many families of Fourier series in
at most \(2d\) integer indices.  Their rational frequencies have height
at most \(D(1+\|\ell\|)\), where \(D\) is quasipolynomial in \(H,B\).
The total absolute coefficient sum is bounded by such a \(D\).

More quantitatively, for any fixed integer \(J\ge1\), one can choose
a quasipolynomial \(D_J\) such that every dyadic shell of indices
\(T\le1+\|\ell\|<2T\) has absolute coefficient sum at most
\[
 D_J T^{-J(d+1)^2}.
\]
For fixed underlying oscillation parameters, the same estimate, after
enlarging \(D_J\), holds for the coefficient envelope obtained by pooling
all rational projection choices of height at most \(B\), separately
for each resulting frequency and effective quadratic part.  The
coefficient envelopes are uniform over amplitudes with the stated
complexity bounds and over the bounded slow values entering their
construction.  In particular the pooled tail can be made at most
\(\varepsilon\) by truncation at a size quasipolynomial in
\(H,B,\varepsilon^{-1}\).
\end{lemma}

\begin{proof}
Put \(u=Vz\), \(u_f=P_fu\), \(u_s=P_su\), and \(w=u-n\).
For the symmetric part,
\[
 u^TC_{\mathrm{sy}}(u-n)
 =\tfrac12u^TC_{\mathrm{sy}}u
  +\tfrac12w^TC_{\mathrm{sy}}w
  -\tfrac12n^TC_{\mathrm{sy}}n.
\]
For the skew part, isotropy on \(U\) gives
\[
 u^TC_{\mathrm{sk}}w
 =u_f^TC_{\mathrm{sk}}u_s
   -u_f^TC_{\mathrm{sk}}P_sn
   +u_s^TC_{\mathrm{sk}}w.
\]
The first term is precisely the extra quadratic included in
\(\Lambda_*\).

Choose an even integer \(P_0\) clearing the denominators of
\(C_{\mathrm{sy}}\).  Then
\(-n^TC_{\mathrm{sy}}n/2\bmod\Z\) is determined by
\(n\bmod P_0\Z^d\).
Also choose a common positive denominator \(q_0\) for \(P_s\).
On the box under consideration, \(u_s=V_sz\) is bounded, and
\(w\) belongs to the support of \(\chi\).
Thus \(P_sn=u_s-P_sw\) ranges over a bounded subset of
\(q_0^{-1}\Z^d\).  The size of this set, \(P_0^d\), and
\(q_0\) are all quasipolynomial in \(H,B\).

Split the sum over \(n\) according to
\[
 P_sn=N,\qquad n\equiv a\pmod{P_0\Z^d}.
\]
For each pair \((N,a)\), remove the fixed residue phase and the rational
linear phase
\[
 -u_f^TC_{\mathrm{sk}}N.
\]
What remains is a smooth amplitude in the two vector variables
\((u_s,w)\).  To see explicitly that it admits a periodic Fourier
expansion, let \(\eta(s)\) be a smooth cutoff equal to one on all
values of \(V_sz\) in the box.  Choose a smooth equality cutoff \(E\)
with \(E(0)=1\), supported in the coordinate cube of radius
\(1/(3q_0)\).  Then the function
\[
 \eta(s)\!
 \sum_{n\in a+P_0\Z^d}
 \chi(u-n)E\bigl(s-P_s(u-n)-N\bigr)
 e\!\left(\tfrac12(u-n)^TC_{\mathrm{sy}}(u-n)
                      +s^TC_{\mathrm{sk}}(u-n)\right)
\]
has the required value when \(s=P_su\).
Indeed its equality-cutoff argument then equals \(P_sn-N\);
the lattice spacing makes the cutoff exactly the indicator of the
specified equality.  It is periodic in \(u\) with period \(P_0\)
in each coordinate.  By taking the support of \(\eta\) strictly inside
a larger integer-period box, it extends smoothly and periodically
in \(s\), with a period \(P_1\) quasipolynomial in \(H,B\).

The resulting Fourier modes have the form
\[
 e\bigl(p^Ts/P_1+\ell^Tu/P_0\bigr),\qquad p,\ell\in\Z^d.
\]
Substitution of \(s=P_su\), and restoration of the rational linear
phase for \(N\), gives
\[
 \sigma=\sigma_N+P_s^Tp/P_1+\ell/P_0,
 \qquad
 \sigma_N^Tu=-(P_fu)^TC_{\mathrm{sk}}N.
\]
These frequencies have the height asserted in the lemma.  The sum
of the two ordinary quadratic terms is
\(z^T\Lambda_*z\), proving the exact identity
\eqref{eq:src-3}.

Here are the quantitative convergence details.  On each of the bounded
supports above, all coefficients in the smooth phases, cutoff scales,
periods, and lattice-point counts are bounded by a quasipolynomial
quantity \(D_0\).  The product rule and the chain rule imply, at order
\(m\), derivative bounds of the form
\[
 D_0^{C(m+d+1)}(m+1)^{C(m+1)} .
\]
This estimate also bounds derivatives after rescaling the periodic
variables to unit periods.  It follows either by differentiating the
quadratic exponential, whose \(m\)-th derivative is a polynomial of
degree at most \(m\) times that exponential, or directly by induction
on \(m\).  The cutoffs satisfy the same fixed-order Gevrey estimate.

For a nonzero integer Fourier index \(v=(p,\ell)\), select a coordinate
with \(|v_j|\ge\|v\|/(2d)^{1/2}\) and integrate by parts \(m\) times
in that coordinate.  Thus
\[
 |\widehat A(v)|
 \le D_0^{C(m+d+1)}
       (m+1)^{C(m+1)}(2d)^{m/2}(1+\|v\|)^{-m},
\]
after increasing the constant to include the zero index.
A dyadic shell has at most \((CT)^{2d}\) indices.
Taking \(m\ge J(d+1)^2+4d+4\) gives the displayed shell bound.
Since \(d\le\log H\), the prefactor for this derivative order is
quasipolynomial in \(H,B\).  Summing the shells proves absolute
convergence and the bound for the full coefficient sum.

There are at most \(B^{C(d+1)^2}\) rational projection matrices
of height at most \(B\).
For each such matrix there are quasipolynomially many choices
\((N,a)\).  All derivative estimates are uniform in the bounded
slow parameters and in amplitudes with the stated complexity bounds.
Consequently we may replace every coefficient by its uniform envelope
before summing over this finite list.  If different indices give the
same frequency, use the sum of their envelopes at that frequency.
Choosing the derivative order larger by \(C(d+1)^2\) absorbs, as well,
the count \((D_0T)^{C(d+1)}\) of bounded-height rational frequencies
at an index scale \(T\).
This proves the pooling assertion in the form stated.
Finally, the sum of the shells above a cutoff \(T_0\) is at most
\(C D_JT_0^{-J(d+1)^2}\); choosing this to be at most
\(\varepsilon\) gives the required quasipolynomial truncation.
\end{proof}

\subsection{Separating two atoms along a progression}

Recall that \(y_j=x+jt\), \(0\le j\le3\).

\begin{lemma}[Progression separation]\label{lem:ap-separation}
Let \(\varphi_i,\varphi_k\) be bounded atoms of complexity at most \(H\)
in two distinct slots \(i,k\), and let \(a,b\) denote the remaining
slots.  Their product has an exact absolutely convergent expansion
in terms
\begin{equation*}
 c\,\psi_a(y_a)\psi_b(y_b)Q(t).
 \tag{4}\label{eq:src-4}
\end{equation*}
All three factors are bounded atoms independent of the interval scale.
For some \(D\) quasipolynomial in \(H\), the terms can be arranged in
dyadic shells \(T\ge1\) with total coefficient mass at most
\(DT^{-100}\).  In a shell \(T\), each factor has dimension at most
\(\log D\), and complexity at most \((DT)^C\); unrestricted real linear
phases are permitted as in Definition~\ref{def:quadratic-atom}.
\end{lemma}

\begin{proof}
We first separate only \(\varphi_i\).  Temporarily suppress its linear
phase, put \(\alpha_i=1\), \(\alpha_k=0\), and choose
\(\alpha_a,\alpha_b\) so that
\[
 \sum_{j=0}^3\alpha_j=0,\qquad
 \sum_{j=0}^3j\alpha_j=0.
\]
The four indices lie in a fixed set, so these coefficients are fixed
rationals of bounded height.

Choose a nonnegative, compactly supported Gevrey function
\(\rho\) on \(\R^d\) whose integer translates sum to one.
Insert its partitions of unity at \(\theta y_a,\theta y_b,\theta t\).
Write the corresponding lifts as
\[
 d_a=\theta y_a-m_a,\qquad
 d_b=\theta y_b-m_b,\qquad
 d=\theta t-m.
\]
For the original summation index \(m_i\), put
\[
 d_i=\theta y_i-m_i,\qquad
 k_j=d_j-d_a-(j-a)d\quad(j=i,b),\qquad k_a=0.
\]
Each \(k_j\) is integral because \(y_j=y_a+(j-a)t\).
The support restrictions bound the possible \(k_i,k_b\)
quasipolynomially in \(H\).
Conversely, fixing these integers determines \(m_i\) from
\(m_a,m\), and the condition for \(k_b\) is enforced by a smooth
equality cutoff in \(d_b-d_a-(b-a)d-k_b\).
Such a cutoff can equal one at zero and vanish at every other integer
vector, so this insertion is exact.

Put \(S=\sum_jj^2\alpha_j\).  The two relations defining the
\(\alpha_j\)'s imply the exact identity
\[
 \sum_j\alpha_j
       \bigl(\lambda y_j^2+(\theta y_j)^TCd_j\bigr)
 =S\bigl(\lambda t^2+(\theta t)^TCd\bigr)
       +\sum_j\alpha_j(\theta y_j)^TCk_j .
\]
Indeed the coefficients of the common lift \(d_a\) and of \(x\)
cancel, and the coefficient of \(t\) in
\(\sum_j(j-a)\alpha_jy_j\) is \(S\).
Solving this identity for the \(i\)-phase leaves the \(a\)- and
\(b\)-phases with coefficients \(-\alpha_a,-\alpha_b\), and the
\(t\)-phase with coefficient \(S\).
The last displayed sum is linear in \(x,t\), and hence is a sum
of a linear function of \(y_a\) and a linear function of \(y_b\).
Restore the suppressed \(\beta y_i\) in the same way.

The remaining coupled amplitude is a smooth function of
\((d_a,d_b,d)\) on a bounded box.  It is a product of the partition
bumps, the equality cutoff, and
\(\chi(d_a+(i-a)d+k_i)\).
Its derivative bounds are quasipolynomial versions of the fixed
Gevrey bounds.  Choose slightly larger product bumps equal to one
on its support, extend the amplitude periodically inside that larger
box, and take its Fourier series.  Every Fourier term factors into
a function of \(d_a\), a function of \(d_b\), and a function of \(d\).
After the independent lattice sums, these are precisely three atoms
of the form \eqref{eq:src-1}, with rational chart matrices
\(-\alpha_aC,-\alpha_bC,SC\).
A Fourier index of size \(T\) contributes only derivatives of size
\(T^m\) at order \(m\) to their smooth amplitudes, so its complexity
is bounded by \((DT)^C\).

Apply the same construction to \(\varphi_k\), now setting the
coefficient of the \(i\)-slot to zero, and combine the two triples
by direct sums of horizontal coordinates.
The number of integral choices and the dimensions are
quasipolynomially bounded in the original complexity.
The integration-by-parts calculation in the proof of
Lemma~\ref{lem:chart-expansion}, with derivative order larger than
the Fourier dimension by at least \(110\), bounds shell coefficient
sums by \(DT^{-100}\).
The number of Fourier indices and integral choices in each bounded
shell is also quasipolynomial.
Normalize the three factors to have supremum norm one and put their
bounded normalization factors into \(c\).  The bounds remain of the
stated form.  Absolute convergence justifies all rearrangements and
proves the identity for every \((x,t)\).
\end{proof}

\begin{lemma}[Single-scale separation]\label{lem:single-scale-separation}
Fix a term in \eqref{eq:src-4}, without its coefficient, of complexity
at most \(H\).  Let \(\zeta\) be a fixed nonnegative smooth function,
equal to one on the range of \(t/r\) allowed by a local kernel \(w_I\),
and supported on a larger fixed interval.  Put
\[
 A_r=\sup_{\xi\in\R}
       \left|\int_\R Q(ru)\zeta(u)e(-\xi u)\,du\right|,
 \qquad r=|I|.
\]
Its bilinear form norm in the remaining two \(L^2(I)\) slots, with
normalized interval norms, is at most \(C A_r\).

If \(A_r\ge\delta>0\), then, for every \(0<\varepsilon<1/2\), the
kernel of this bilinear form, extended by zero off its support, can be
approximated uniformly on \(I\times I\), to error \(\varepsilon\),
by a finite separated sum
\(\sum_\nu c_\nu\phi_{\nu,a}(y_a)\phi_{\nu,b}(y_b)\).
The \(\phi_{\nu,e}\) are bounded atoms, and their complexities, their
number, and \(\sum_\nu|c_\nu|\) are quasipolynomial in
\(H,\delta^{-1},\varepsilon^{-1}\).

Before separating the smooth kernel, the retained terms have phases
\(\lambda_*t^2+\gamma t\), where
\(|\lambda_*|r^2\) is quasipolynomially bounded.
The quantities \(\lambda_*\) are bounded-height rational combinations
of \(\lambda\) and the entries of \(\theta\otimes\theta\); each
\(\gamma-\beta\) is a bounded-height rational combination of \(\theta\).
The height bounds are quasipolynomial in the same parameters.
\end{lemma}

\begin{proof}
Use \((y_a,y_b)\) as coordinates; their Jacobian is \(|b-a|\) and
\(t=(y_b-y_a)/(b-a)\).
After normalization by \(r\), the smooth kernel has a Fourier
expansion on a fixed two-dimensional box whose coefficient sum is
uniformly bounded.  Each Fourier term merely modulates the two
inputs.  Insert \(\zeta(t/r)\), which is one on the relevant support.
The remaining operator is convolution with
\(Q(t)\zeta(t/r)\), after the fixed change of scale \(b-a\).
Its Fourier multiplier is bounded by \(C r A_r\).
Plancherel, the factor \(r^{-2}\) in the definition of \(H_I\), and
\(\|f\|_{L^2(I)}=r^{1/2}\|f\|_{2,I}\) prove the first assertion.
The factors \(\psi_a,\psi_b\) are bounded multipliers and do not
increase the two \(L^2\) norms.

For the second assertion apply Lemma~\ref{lem:weyl} to \(Q(r\,\cdot)\),
with \(V=r\theta\) and \(\Lambda=r^2\lambda\).
Lemma~\ref{lem:chart-expansion}, truncated to a sufficiently small
uniform error, then gives terms
\[
 e(\lambda_*t^2+\gamma t),\qquad
 \gamma=\beta+\sigma^T\theta,\qquad
 |\lambda_*|r^2\le D .
\]
The formula for \(\Lambda_*\) in \eqref{eq:src-2} gives exactly the
asserted rational dependence of \(\lambda_*\).
Split the linear phase by
\[
 e(\gamma t)=
 e\bigl(-\gamma y_a/(b-a)\bigr)
 e\bigl(\gamma y_b/(b-a)\bigr).
\]
The remaining smooth two-variable kernel,
\(w_I(x,t)e(\lambda_*t^2)\), has derivatives in normalized
coordinates bounded by
\(D^{C(m+1)}(m+1)^{C(m+1)}\) at order \(m\).
Its Fourier expansion, with fixed product cutoffs, has
quasipolynomial coefficient sum and a quasipolynomial truncation
giving the prescribed error.
The resulting factors are a bounded one-variable smooth bump at
scale \(r\), a linear phase, and \(\psi_e\).

For completeness, a smooth bump on \(I\), restricted to \(I\), is
itself represented by a chart with zero matrix and one additional
slow horizontal coordinate.  Map \(I\) to an interval strictly
inside one lattice cell, put the bump in the corresponding amplitude,
and support that amplitude away from every other lattice translate.
On \(I\) only the chosen translate contributes.
Lemma~\ref{lem:chart-operations} handles the affine offset and the
product with \(\psi_e\).  Outside \(I\) a bounded periodic extension
is harmless.  This realizes all separated factors as bounded atoms.
Choose the successive truncation errors divided by the already
bounded coefficient sums.  Only a fixed number of quasipolynomial
budgets is composed, giving all the stated estimates.
\end{proof}

\subsection{A bounded continuous inverse theorem}

We use the following external input, only at degree two.

\begin{theorem}[Leng--Sah--Sawhney]\label{thm:lss}
For a \(1\)-bounded function \(f:[L]\to\C\) with
\(\|f\|_{U^3[L]}\ge\delta\), there are a degree-two filtered
nilmanifold \(G/\Gamma\), a polynomial sequence \(g:\Z\to G\)
adapted to its filtration, and a function \(F:G/\Gamma\to\C\) such that
\[
 \left|\frac1L\sum_{n=1}^L
        f(n)\overline{F(g(n)\Gamma)}\right|\ge\varepsilon.
\]
The dimension is at most \(C(1+\log(1/\delta))^{C}\), and the complexity
of \(G/\Gamma\), the Lipschitz norm of \(F\), and
\(\varepsilon^{-1}\) are at most
\(\exp(C(1+\log(1/\delta))^C)\).
The nilmanifold has a rational Malcev basis adapted to the filtration,
with integer second-kind coordinates for \(\Gamma\).
The Lipschitz norm includes the supremum norm.
\end{theorem}

This is the \(s=2\) instance of \cite[Theorem 1.2]{LSS2024};
the coordinate and complexity conventions are
\cite[Definitions 2.3--2.4 and 3.3--3.5]{LSS2024}.
In particular no mean-zero hypothesis is imposed.

For the continuous version, first identify an interval with \([0,1)\).
Fix the circle \(\T_{16}=\R/(16\Z)\), with its probability Haar measure,
and extend functions on \([0,1)\) by zero to this circle.  Define
\[
 \|g\|_{\mathcal U}^8
   =\int_{\T_{16}^4}
       \prod_{\omega\in\{0,1\}^3}
       \mathcal C^{|\omega|}g(x+\omega\cdot h)
                 \,d\mu(x)\,d\mu(h_1)\,d\mu(h_2)\,d\mu(h_3),
\]
where \(\mathcal C\) denotes complex conjugation.
On an arbitrary interval \(I\), \(\mathcal U(I)\) denotes this
norm after the affine normalization to \([0,1)\).
All constants caused by this fixed circle length are absolute.

\begin{lemma}[Elementary cube bounds]\label{lem:cube-bounds}
The quantity \(\|\cdot\|_{\mathcal U}\) is a seminorm.
The mixed cube form with possibly different functions \(g_\omega\)
at its eight vertices satisfies
\[
 \left|\int_{\T_{16}^4}
       \prod_\omega\mathcal C^{|\omega|}
            g_\omega(x+\omega\cdot h)\,d\mu(x)\,d\mu(h)\right|
 \le \prod_\omega\|g_\omega\|_{L^2(\T_{16})}.
\]
Consequently cube forms are continuous in \(L^2\), and
\(\|g\|_{\mathcal U}\le C\|g\|_{2,[0,1)}\) for supported functions.
\end{lemma}

\begin{proof}
The usual cube Cauchy--Schwarz inequality is valid on this circle.
One direct derivation writes a vertex as
\(x_1^{\omega_1}+x_2^{\omega_2}+x_3^{\omega_3}\), using three
independent pairs of Haar variables.
At the first Cauchy--Schwarz step the factors with the first
coordinate equal to zero and one are put on opposite sides; repeating
for the second and third coordinates gives
\[
 |\Lambda_{\mathrm{cube}}(g_\omega)|
       \le\prod_\omega\|g_\omega\|_{\mathcal U}.
\]
The same iterations for a single function express its eighth power
as an integral of squared absolute values, proving nonnegativity.
Expanding the cube of \(f+g\), applying the displayed inequality to
each mixed term, and summing the binomial coefficients gives the
triangle inequality.

For the stated \(L^2\) estimate, split the vertices into
\[
 \{000,100,010,001\}
       \quad\hbox{and}\quad
 \{111,011,101,110\}.
\]
Each four-tuple of corresponding linear forms is an integer linear
automorphism of the four Haar variables, with determinant of absolute
value one.  Apply Cauchy--Schwarz to the product of these two
four-factor groups, then change variables in each squared integral.
Each squared integral factors into the four squared \(L^2\) norms.
This proves the estimate.  Expanding the difference of two mixed cube
forms proves their continuity, first for bounded functions and then
for all \(L^2\) functions by truncation.  The fixed normalization
between circle measure and interval measure proves the final assertion.
\end{proof}

\begin{lemma}[Continuous bounded inverse]\label{lem:continuous-inverse}
For \(0<\delta<1/2\), let \(g\) be supported on an interval \(I\),
with \(\|g\|_\infty\le1\) and
\(\|g\|_{\mathcal U(I)}\ge\delta\).
There is a bounded atom \(\varphi\) of complexity at most
\[
 H_\delta=\exp\!\left(C(1+\log(1/\delta))^C\right)
\]
such that
\[
 \left|\fint_I g(y)\overline{\varphi(y)}\,dy\right|
       \ge H_\delta^{-1}.
\]
\end{lemma}

\begin{proof}
Normalize \(I=[0,1)\).
For \(L\ge1\), define the cell averages
\[
 g_L(j)=L\int_{j/L}^{(j+1)/L}g(y)\,dy,\qquad 0\le j<L,
\]
and extend them by zero to \(\Z/(16L\Z)\).
Their discrete raw cube averages converge to the continuous raw
cube average above.  Here is a justification that needs no regularity
of \(g\).
For a continuous bounded function the assertion is Riemann-sum
convergence, because its cell averages differ from endpoint values
by its modulus of continuity.
For two functions bounded by one, telescoping the eight-factor product
bounds the difference of their continuous cube averages by
\(8\|g-b\|_{L^1(\T_{16})}\).
The identical discrete bound holds for their cell averages, since
\[
 \frac1L\sum_{j=0}^{L-1}|g_L(j)-b_L(j)|
       \le\|g-b\|_{L^1([0,1))}.
\]
Approximation in \(L^1\) by continuous functions bounded by one proves
the assertion.

There is no wraparound contribution hidden in this comparison.
For a nonzero supported cube choose \(x\in[0,1)\), and lift
\(h_j\) as the difference between \(x+h_j\in[0,1)\) and \(x\).
Then \(|h_j|\le1\), and every real vertex belongs to \([-3,4]\).
Within that range, reduction modulo \(16\) lies in \([0,1)\) exactly
when the real vertex does.
Thus the circle and interval zero-extension cube averages differ
only by their fixed normalizing factors.  The denominator defining
\(U^3[L]\), namely the cube norm of the interval indicator, likewise
converges to a fixed positive constant.  For all sufficiently large
\(L\), therefore,
\[
 \|g_L\|_{U^3[L]}\ge c\delta,
\]
after shifting the integer labels by one if needed.
No bound on this choice of \(L\) is required.

Apply Theorem~\ref{thm:lss}.  Write its dimension as \(d_0\), its
complexity as \(M_0\), its Lipschitz bound as \(K_0\), and its
correlation as \(\varepsilon_0\).
All these quantities have the bounds asserted there.
We next convert its nilsequence into chart atoms, tracking only
quantities independent of the real orbit coefficients.

In second-kind Malcev coordinates adapted to the central filtration
subgroup \(G_2\), write a group element as \((z,s)\), with
\(z\in\R^p\), \(s\in\R^{d_0-p}\).
Since \([G,G]\subset G_2\) and \([G,G_2]=1\), commuting the
one-parameter coordinate factors into their prescribed order gives
the group law
\[
 (z,s)(z',s')=(z+z',s+s'+B(z,z')).
\]
The coefficients of the bilinear map \(B\) are rational and have
height bounded quasipolynomially in \(M_0\) and \(e^{d_0}\).
Indeed, in step two every interchange contributes just its central
commutator with coefficient \(z_i z'_j\); all further commutators
vanish.  The rational structure constants in the Malcev basis are
part of the nilmanifold complexity.  Integer coordinate tuples
represent exactly the lattice \(\Gamma\).

Write the polynomial sequence in these coordinates as
\((z(n),s(n))\).
Its second group differences have no horizontal part, so
\(z(n)=z_0+n\theta\).
For the unit right difference \(Dg(n)=g(n+1)g(n)^{-1}\), the group law
gives
\[
 Dg(n)=\bigl(\theta,\Delta s(n)-B(\theta,z(n))\bigr),\quad
 D^2g(n)=\bigl(0,\Delta^2s(n)-B(\theta,\theta)\bigr),\quad
 D^3g(n)=\bigl(0,\Delta^3s(n)\bigr).
\]
The degree-two filtration makes \(D^3g=1\).  Thus \(s(n)\) is an
ordinary vector polynomial of degree at most two.
No bound is imposed on the coefficients of \(z(n),s(n)\).

Choose \(\rho\) as the product of fixed one-dimensional nonnegative
Gevrey partitions supported in \((-1,1)\), so that
\(\sum_{m\in\Z^p}\rho(z-m)=1\).
Right multiplication by the lattice element \((-m,0)\) gives
\[
 F(g(n)\Gamma)
 =\sum_{m\in\Z^p}\rho(z(n)-m)
       \widetilde F\bigl(z(n)-m,s(n)-B(z(n),m)\bigr),
 \qquad
 \widetilde F(w,u)=F((w,u)\Gamma).
\]
The function \(\widetilde F\) is periodic in every central coordinate.
On a bounded horizontal box its coordinate Lipschitz constant is
quasipolynomially bounded in \(K_0,M_0,e^{d_0}\).
To check this directly, the coordinate difference of two nearby group
elements is
\[
 (w,u)(w',u')^{-1}
   =\bigl(w-w',\,u-u'-B(w-w',w')\bigr);
\]
on a bounded box its coordinates are bounded by a controlled multiple
of the coordinate displacement.  The Malcev metric is bounded by
the length of this coordinate displacement path.  This proves the
claimed coordinate Lipschitz estimate.

Multiply by a fixed horizontal cutoff equal to one on \(\supp\rho\)
and supported strictly inside a larger period box.
Periodize horizontally, and convolve in all coordinates with product
Fej\'er kernels of order \(N_0\).
For a Lipschitz function on a circle the error is bounded by its
Lipschitz constant times \(C\log(2N_0)/N_0\), since this is an upper
bound for the kernel's first moment.  Telescoping over the \(d_0\)
coordinates gives the same estimate multiplied by \(d_0\).
Choose \(N_0\) so that this error is at most
\(\varepsilon_0/4\).  It may be chosen quasipolynomially in
\(K_0,M_0,e^{d_0},\varepsilon_0^{-1}\).
The resulting trigonometric polynomial has the form
\[
 \sum_{\ell,k}a_{\ell,k}
             e(\ell^Tw/P_2+k^Tu),
 \qquad \|\ell\|_\infty,\|k\|_\infty<N_0,
\]
where \(P_2\) is a fixed integer horizontal period.
Each coefficient has absolute value at most a controlled multiple
of \(K_0\), so its absolute coefficient sum is at most
\(C K_0(2N_0+1)^{d_0}\).
Because the partition \(\rho\) is nonnegative and sums to one,
the approximation error remains at most \(\varepsilon_0/4\)
after summing over \(m\).

For a fixed central frequency \(k\), define \(C_k\) by
\(z^TC_kw=k^TB(z,w)\).
The associated summand of this approximation is
\[
 e\bigl(k^Ts(n)-z(n)^TC_kz(n)\bigr)
 \sum_{m\in\Z^p}
     \rho(z(n)-m)e\bigl(\ell^T(z(n)-m)/P_2\bigr)
     e\bigl(z(n)^TC_k(z(n)-m)\bigr).
\]
Its outside phase is an ordinary quadratic in \(n\).
The matrix \(C_k\) is rational with controlled height, while
the horizontal orbit is affine.
Lemma~\ref{lem:chart-operations} removes its constant offset,
without putting an unbounded parameter in the amplitude.
Thus each summand is an atom \eqref{eq:src-1}.
The positivity of \(\rho\) shows that it has supremum at most one
before the harmless normalization of the coefficients.
Its complexity and the total absolute coefficient sum are
quasipolynomial in \(\delta^{-1}\).
It follows that one such discrete atom \(a(n)\) satisfies
\[
 \left|\frac1L\sum_{j=0}^{L-1}
                g_L(j)\overline{a(j)}\right|\ge\kappa,
 \qquad
 \kappa^{-1}\le
 \exp\!\left(C(1+\log(1/\delta))^C\right).
\]
Here and below a shift of the polynomial argument incorporates the
choice of integer labels.

We finish with an exact floor interpolation, which is necessary because
the chosen atom may oscillate on the sampling scale.
Write its parameters as \((\lambda,\beta,\theta,C,\chi)\).
Let \(P\) clear the denominators of \(C\), and write
\(\theta=\theta_0+Pq\), with \(q\in\Z^d\) and
\(0\le(\theta_0)_j<P\).
At integer arguments \(n\), reindexing the chart sum by \(Pqn\)
replaces the ordinary quadratic coefficient by
\(\lambda+(Pq)^TC\theta_0\); the other extra term is integral.
Reduce this new coefficient and \(\beta\) modulo one.
The same sampled atom consequently has a representation with
\[
 0\le\lambda,\beta<1,\qquad
 0\le\theta_j<P.
\]
Only this equality at integer arguments is needed.

Put \(\alpha=2\lambda+\theta^TC\theta\), and define the enlarged
horizontal vector and matrix
\[
 \Theta=\begin{pmatrix}\theta\\1\\\alpha\end{pmatrix},
 \qquad
 \widetilde C=
 \begin{pmatrix}
 C&0&0\\
 0&0&0\\
 0&-1&0
 \end{pmatrix}.
\]
Choose \(0\le\eta\le1\), Gevrey of the fixed order at cutoff scale
\(\epsilon\), supported in \((0,1)\), and equal to one outside
endpoint strips of total length at most \(2\epsilon\).
Choose a one-dimensional Gevrey partition
\(\sum_{\ell\in\Z}\rho_0(r-\ell)=1\).
For lift variables \((w,h,r)\), define
\[
 \widetilde\chi(w,h,r)=
   \chi(w-\theta h)\eta(h)\rho_0(r)
   e\!\left(-h\theta^TCw
             +(\lambda+\theta^TC\theta)h^2-\beta h\right).
\]
The atom
\[
 A(v)=e(\lambda v^2+\beta v)
       \sum_{n,m,\ell}
       \widetilde\chi(\theta v-n,v-m,\alpha v-\ell)
       e\!\left((\Theta v)^T\widetilde C
                        (\Theta v-(n,m,\ell))\right)
\]
satisfies the exact identity
\[
 A(v)=\eta(\{v\})a(\lfloor v\rfloor)
 \qquad(v\in\R),
\]
with value zero at integer \(v\).
Indeed the support of \(\eta\) forces \(m=\lfloor v\rfloor\);
the \(\ell\)-sum is one.  With \(h=v-m\) and \(w=\theta v-n\),
the phase in the remaining summand is
\[
 \lambda v^2+\beta v+(\theta v)^TCw
 -(2\lambda+\theta^TC\theta)vh-h\theta^TCw
       +(\lambda+\theta^TC\theta)h^2-\beta h,
\]
which equals
\(\lambda m^2+\beta m+(\theta m)^TC(\theta m-n)\).
This proves the identity.

All parameters occurring in the new amplitude are now bounded:
the reductions above bound \(\lambda,\beta,\theta\) and hence
\(\alpha\), while its support bounds \(h,r,w-\theta h\).
The cutoff \(\eta\) has Gevrey derivative bounds with scale
\(\epsilon^{-1}\).
Thus the complexity of \(A\) is quasipolynomial in the old complexity
and \(\epsilon^{-1}\); its two extra horizontal coordinates and its
new matrix entry \(-1\) introduce no uncontrolled height.
Moreover \(|A|\le1\), because its values are a cutoff times sampled
values of the bounded atom.

Take \(v=Ly\).  This changes only unrestricted real oscillation
parameters.  Exact cell averaging gives
\[
 \int_0^1 g(y)\overline{a(\lfloor Ly\rfloor)}\,dy
       =\frac1L\sum_{j=0}^{L-1}g_L(j)\overline{a(j)}.
\]
Replacing the sampled atom by \(A(Ly)\) changes this integral by at
most \(2\epsilon\), since the total measure of the removed cell
endpoint strips is at most \(2\epsilon\).
Choose \(\epsilon\le\kappa/8\).
The resulting correlation is at least \(3\kappa/4\), and its
complexity is quasipolynomial in \(\delta^{-1}\), independently of
\(L\).  This proves the normalized-interval assertion.
Lemma~\ref{lem:chart-operations} transfers it to the original interval.
\end{proof}

\subsection{A bounded-input decomposition}

The argument below uses the inverse-to-three-part decomposition method of
Gowers--Wolf \cite[Proposition~3.3 and Theorem~3.4]{GW2010}, adapted here to
the local interval and the \(L^2\)-closure step.  The quantitative inverse
input remains the Leng--Sah--Sawhney theorem \cite{LSS2024} used in
Lemma~\ref{lem:continuous-inverse}.

\begin{lemma}[Structured, uniform, and \(L^1\) decomposition]
\label{lem:bounded-decomposition}
For \(0<\tau<1/2\) and a function \(f\) on \(I\) with
\(\|f\|_\infty\le1\), there is a finite decomposition
\[
 f=\sum_\nu c_\nu\varphi_\nu+u+v
 \quad\hbox{on }I
\]
such that the \(\varphi_\nu\) are bounded atoms of complexity at most
\(D_\tau\),
\[
 \sum_\nu|c_\nu|\le D_\tau,\qquad
 \|u\|_{\mathcal U(I)}\le\tau,\qquad
 \|v\|_{1,I}\le\tau,
\]
and \(u,v\in L^2(I)\).  Here
\(D_\tau\le\exp(C(1+\log(1/\tau))^C)\).
\end{lemma}

\begin{proof}
All functions and inner products in this proof use normalized measure
on \(I\), and all functions are extended by zero for their
\(\mathcal U(I)\) norms.
Put \(\sigma=\tau/2\).
Let \(\mathcal A\) be the bounded atoms of the complexity supplied by
Lemma~\ref{lem:continuous-inverse} at threshold \(\sigma^2\), restricted
to \(I\), and let \(\kappa>0\) be the corresponding correlation bound.
Choose \(M\ge\max(1,2\sigma/\kappa)\), quasipolynomial in
\(\sigma^{-1}\).  Consider the convex balanced set
\[
 \mathcal B=
 M\operatorname{absconv}(\mathcal A)
 +\{u\in L^2(I):\|u\|_{\mathcal U(I)}\le\sigma\}
 +\{v\in L^2(I):\|v\|_{1,I}\le\sigma\}.
\]
Here the absolute convex hull consists of finite combinations.
This set contains an \(L^2\) ball about zero.

Suppose \(f\) were outside the \(L^2\) closure of \(\mathcal B\).
The real Hahn--Banach separation theorem, and the complex Riesz
representation of a real functional as a real part, give
\(h\in L^2(I)\), normalized so that
\[
 \operatorname{Re}\langle f,h\rangle>1,\qquad
 \sup_{z\in\mathcal B}\operatorname{Re}\langle z,h\rangle\le1.
\]
Balancedness of each summand then implies
\[
 \|h\|_{\infty,I}\le\sigma^{-1},\qquad
 |\langle u,h\rangle|
       \le\sigma^{-1}\|u\|_{\mathcal U(I)}
          \quad(u\in L^2(I)),\qquad
 \sup_{\varphi\in\mathcal A}|\langle h,\varphi\rangle|\le M^{-1}.
\]
The first inequality follows by testing bounded simple functions in
the \(L^1\) ball; no assumption that the \(L^1\) dual is represented
by an arbitrary measure is needed.
Since \(\|f\|_{2,I}\le1\), the first strict pairing inequality gives
\(\|h\|_{2,I}>1\).
Testing the second bound with \(u=h\) yields
\[
 \|h\|_{\mathcal U(I)}
       \ge\sigma\|h\|_{2,I}^2>\sigma.
\]
Thus \(\sigma h\) is bounded by one and has
\(\mathcal U(I)\) norm greater than \(\sigma^2\).
Lemma~\ref{lem:continuous-inverse} supplies \(\varphi\in\mathcal A\)
with
\(|\langle h,\varphi\rangle|\ge\kappa/\sigma>M^{-1}\),
a contradiction.

Therefore \(f\) belongs to the closure of \(\mathcal B\).
Choose an element of \(\mathcal B\) within \(\tau/2\) in \(L^2(I)\).
Its finite structured combination has coefficient sum at most \(M\).
Put the approximation error into its \(L^1\) error term, using
\(\|\cdot\|_{1,I}\le\|\cdot\|_{2,I}\).
The \(L^1\) error is then at most \(\tau\), and the uniformity error is
at most \(\tau/2\).  All components lie in \(L^2(I)\), as required.
\end{proof}

\subsection{Detection from a local progression form}

We first state explicitly the form estimate used with the uniform part
of the preceding decomposition.

\begin{lemma}[A generalized von Neumann bound with one \(L^2\) input]
\label{lem:local-von-neumann}
For four distinct slots occupied by \(R,g,b,c\), one has
\[
 |H_I(R,g,b,c)|
   \le C\|R\|_{2,I}\|g\|_{\mathcal U(I)}
                      \|b\|_{\infty,I}\|c\|_{\infty,I}.
\]
One also has
\[
 |H_I(R,g,b,c)|
   \le C\|R\|_{1,I}\|g\|_{1,I}
                      \|b\|_{\infty,I}\|c\|_{\infty,I}
\]
for \(g\in L^1(I)\), with the first estimate interpreted for
\(R,g\in L^2(I)\).
\end{lemma}

\begin{proof}
Normalize \(I=[0,1)\).
Extend the kernel smoothly to a sufficiently large fixed periodic
\((x,t)\) box, supported where its real arguments lie in \(I\).
The interior support margin permits this extension without changing
the kernel on its support.
Its Fourier coefficient sum is bounded by a constant determined by
the fixed kernel bounds.  Thus it suffices to prove the first estimate
for each character-weighted progression average on the fixed circle.

Let \(R\) occupy slot \(j\), \(g\) slot \(i\), and the two bounded
inputs slots \(k,\ell\).
Disintegrate Haar measure along \(y_j=x+jt\) and apply
Cauchy--Schwarz in that variable.
The first cost is \(\|R\|_2^2\), and translation along the free
direction \((-j,1)\) introduces an independent Haar difference
parameter \(h_1\) in every remaining factor.
Next disintegrate along \(y_k\), keeping \(h_1\) on the base.
The repeated \(b\) factors are unchanged under translations in
\((-k,1)\); Cauchy--Schwarz therefore removes them at cost
\(\|b\|_\infty^4\) and introduces \(h_2\).
The corresponding step at slot \(\ell\) removes the four repeated
\(c\) factors at cost \(\|c\|_\infty^8\) and introduces \(h_3\).
After these three steps the eighth power of the original average
is bounded by
\[
 \|R\|_2^8\|b\|_\infty^8\|c\|_\infty^8
 \int\prod_{\omega\in\{0,1\}^3}
    \mathcal C^{|\omega|}
    g\bigl(y_i+\omega_1(i-j)h_1
                  +\omega_2(i-k)h_2+\omega_3(i-\ell)h_3\bigr).
\]
Every coefficient \(i-j,i-k,i-\ell\) is a nonzero integer.
Multiplication by it preserves Haar probability measure, and the
three parameters remain independent.  The last integral is exactly
\(\|g\|_{\mathcal U}^8\).

A Fourier character of \((x,t)\) may be distributed among slots
\(0,1\), since \(x=y_0\), \(t=y_1-y_0\).
Multiplication by a circle character preserves the \(L^2\) and
supremum norms and the \(U^3\) norm, so the same estimate applies to
every Fourier term.  The argument starts with bounded inputs; the
mixed-cube estimate in Lemma~\ref{lem:cube-bounds} and \(L^2\)
approximation extend it to the stated class.
Summing the kernel's absolutely convergent Fourier series proves
the first bound.

For the second bound, use the two distinct coordinates occupied by
\(R\) and \(g\).  Their Jacobian is a fixed nonzero integer.
Bound the other two inputs and the kernel in supremum norm.
The double integral then factors into the two \(L^1\) norms,
giving the assertion after normalization.
\end{proof}

\begin{lemma}[Local detection]\label{lem:detection}
Fix \(q>2\) and \(C_1\ge1\).
Suppose that \(R,h,u,v\) occupy the four slots in any order and satisfy
\[
 \|R\|_{2,I}\le C_1,\qquad
 \|h\|_{q,I}\le1,\qquad
 \|u\|_{2,I},\|v\|_{2,I}\le1,\qquad
 |H_I(R,h,u,v)|>\delta ,
\]
where \(0<\delta<1/2\).
Then \(R\) has a normalized moment against a bounded atom
\(\varphi\) of magnitude at least \(D^{-1}\), where the complexity of
\(\varphi\) is at most \(D\) and
\[
 D\le
 \exp\!\left(C_q(1+\log C_1+\log(1/\delta))^C\right).
\]
The absolute exponent \(C\) is independent of \(q\).
The atom can be conjugated if the desired moment convention is
bilinear rather than sesquilinear.
\end{lemma}

\begin{proof}
We keep track of the finite sequence of tolerances in the argument.
Write \(q'=q/(q-1)<2\), and truncate each opposing input at absolute
value \(B\ge1\), by putting its tail equal to zero.
For the \(h\) tail,
\[
 \|h\1_{\{|h|>B\}}\|_{2,I}\le B^{-(q-2)/2}.
\]
For the \(u,v\) tails, their normalized \(L^2\) bounds give
\[
 \|u\1_{\{|u|>B\}}\|_{q',I},
 \ \|v\1_{\{|v|>B\}}\|_{q',I}
       \le B^{-(q-2)/q}.
\]
Use the local absolute form estimate with exponents
\((2,2,2,2)\) for the first tail and
\((2,q,q',2)\) for either of the other tails.
In each case the reciprocal exponent sum is at most two.
It follows, by replacing the inputs successively, that the total
clipping error is at most \(C C_1B^{-(q-2)/q}\).
Choose
\[
 B=\max\!\left(2,(100C C_1/\delta)^{q/(q-2)}\right).
\]
The clipped form has magnitude at least \(3\delta/4\).
All clipped inputs have supremum norm at most \(B\), and retain
their original \(L^q\) or \(L^2\) bounds.

Apply Lemma~\ref{lem:bounded-decomposition} to the clipped \(h/B\)
with
\[
 \tau_1=\frac{\delta}{100C C_1B^3}.
\]
On multiplying back by \(B\), write its structured part as
\(s_1=\sum_\nu a_\nu\varphi_\nu\), with
\[
 \sum_\nu|a_\nu|\le A_1,\qquad A_1=B D_{\tau_1}.
\]
Lemma~\ref{lem:local-von-neumann} bounds the two error contributions
by \(C C_1B^3\tau_1\) each: the other two opposing inputs are bounded
by \(B\), and \(\|R\|_{1,I}\le C_1\).
The resulting form with \(s_1\) in this slot still has magnitude
greater than \(\delta/2\).

Next apply Lemma~\ref{lem:bounded-decomposition} to the clipped \(u/B\)
with
\[
 \tau_2=\frac{\delta}{100C C_1A_1B^2}.
\]
Its structured part \(s_2=\sum_\mu b_\mu\psi_\mu\) has coefficient sum
\[
 \sum_\mu|b_\mu|\le A_2,\qquad A_2=B D_{\tau_2}.
\]
Sum the error estimates separately over the atoms of \(s_1\).
The two new error contributions are each at most
\(C C_1A_1B^2\tau_2\).
Therefore, denoting the remaining clipped input by \(v_B\),
\[
 |H_I(R,s_1,s_2,v_B)|>\delta/4.
\]
By the absolute coefficient bounds, some bounded atom pair
\(\varphi_i,\varphi_k\) satisfies
\[
 |H_I(R,\varphi_i,\varphi_k,v_B)|
       >\delta_0,\qquad
 \delta_0=\frac{\delta}{4A_1A_2}.
\]
The complexities of this pair and \(\delta_0^{-1}\) are
quasipolynomial in \(C_1,\delta^{-1}\), with a \(q\)-dependent
constant but an absolute log exponent.  Notice also that
\(\|v_B\|_{2,I}\le1\).

Apply Lemma~\ref{lem:ap-separation} to this pair.
Let \(D_0\) bound its absolute coefficient sum and its shell
constants.  For each term without its coefficient the absolute
local form is at most \(C C_1\), since its three atom factors are
bounded and the two free inputs have the indicated \(L^2\) sizes.
Discard shells above \(T_0\), with \(T_0\) chosen so that
\(C C_1D_0T_0^{-100}<\delta_0/4\).
The retained terms have quasipolynomial complexity.
At least one such term has form magnitude at least
\(\delta_1=\delta_0/(2D_0)\).
The first part of Lemma~\ref{lem:single-scale-separation} forces its
Fourier supremum \(A_r\) to be at least
\(\delta_1/(C C_1)\).

Use the second part of that lemma with a uniform kernel error at most
\(\delta_1/(10C C_1)\).
Pair independence bounds the resulting form error by
\(\delta_1/10\).  The surviving form is a finite sum
\[
 \sum_\nu c_\nu
    \left(\fint_I R(y)\phi_{\nu,a}(y)\,dy\right)
    \left(\fint_I v_B(y)\phi_{\nu,b}(y)\,dy\right),
\]
with a quasipolynomial bound \(D_1\) for
\(\sum_\nu|c_\nu|\) and for all atom complexities.
Fixed coordinate Jacobians have been included in the coefficients.
Every second moment has magnitude at most one, by the \(L^2(I)\)
bound on \(v_B\) and the boundedness of its atom.
Consequently one of the first moments has magnitude at least
\(\delta_1/(2D_1)\).
Conjugating its atom if needed gives the stated normalized moment.

There are exactly two bounded-input decompositions, followed by the
progression separation and the single-scale separation.
Thus all budgets above involve a fixed number of quasipolynomial
compositions.  The only \(q\)-dependent exponent used before these
compositions is the power defining \(B\), whose logarithm is a
\(q\)-dependent constant times
\(1+\log C_1+\log(1/\delta)\).
This proves the final quantitative assertion.
\end{proof}
